\begin{ouverture}[calcul]{3}{Calculer avec les décimaux}{Partie 1 · Nombres et calculs}
{Addition · Soustraction · Multiplication · Contrôle}
\objectif{Additionner et soustraire des nombres décimaux.}
\objectif{Multiplier par 10, 100, 1 000 et par 0,1, 0,01, 0,001.}
\objectif{Comprendre, calculer et contrôler le produit de deux décimaux.}
\prerequis{Tables de multiplication ; opérations posées sur les entiers.}
\end{ouverture}
\pagestyle{otMain}\essentiel\fontsize{11}{15}\selectfont
\phantomsection\label{cours-3-1}
\essSection{1}{Additionner et soustraire}
\begin{savoirfaire}[Aligner les rangs]Pour poser une addition ou une soustraction, on aligne les virgules, donc les rangs, et on complète par des zéros : $12{,}70+3{,}85=16{,}55$ ; $20{,}00-7{,}35=12{,}65$. Avant de poser, on estime : $13+4=17$, proche de $16{,}55$.\end{savoirfaire}
\essSection{2}{Multiplier ou diviser par 10, 100, 1 000}
\begin{propriete}Multiplier par 10, 100 ou $1\,000$ rend chaque chiffre 10, 100 ou $1\,000$ fois plus grand : il avance de 1, 2 ou 3 rangs vers la gauche. $4{,}56\times100=456$. Diviser par 10, 100 ou $1\,000$ le fait reculer d’autant : $37{,}2:100=0{,}372$.\end{propriete}
\label{fig-cours-3-1-cours-glisse}\figureGeo{cours-glisse}

\begin{propriete}Multiplier par $0{,}1$, c’est prendre un dixième, donc diviser par 10. De même, multiplier par $0{,}01$ revient à diviser par 100, et par $0{,}001$ à diviser par $1\,000$ : $45\times0{,}1=4{,}5$ ; $6{,}2\times0{,}01=0{,}062$.\end{propriete}
\clearpage
\phantomsection\label{cours-3-2}
\essSection{3}{Multiplier deux nombres décimaux}
\begin{exemple}Un rectangle de $2{,}3$ cm sur $1{,}4$ cm se découpe en quatre parties d’aires $2\times1=2$ ; $2\times0{,}4=0{,}8$ ; $0{,}3\times1=0{,}3$ et $0{,}3\times0{,}4=0{,}12$. Son aire vaut $3{,}22\ \mathrm{cm}^2$ : donc $2{,}3\times1{,}4=3{,}22$.\end{exemple}
\label{fig-cours-3-2-cours-rectangle-produit}\figureGeo{cours-rectangle-produit}

\begin{propriete}$2{,}3=23\times0{,}1$ et $1{,}4=14\times0{,}1$. On peut changer l’ordre des facteurs : $2{,}3\times1{,}4=23\times14\times0{,}1\times0{,}1=322\times0{,}01=3{,}22$. On multiplie sans les virgules ; le résultat a autant de décimales que les deux facteurs réunis.\end{propriete}
\essSection{4}{Organiser et contrôler un calcul}
\begin{savoirfaire}[Ordre de grandeur]$19{,}8\times5{,}1$ est proche de $20\times5=100$. On calcule $198\times51=10\,098$ et on place deux décimales : $100{,}98$, cohérent avec l’ordre de grandeur.\end{savoirfaire}
\begin{attention}Multiplier par un nombre plus petit que 1 donne un résultat plus petit : $8\times0{,}5=4$. Multiplier n’agrandit pas toujours.\end{attention}

\clearpage\rubrique{Découverte}
\begin{blocLivre}{ex-3-1}\begin{exercice}[1]{Additions posées}
Poser et calculer $12{,}7+3{,}85$, puis $45{,}08+9{,}9$.
\end{exercice}\end{blocLivre}
\begin{blocLivre}{ex-3-2}\begin{exercice}[1]{Soustractions posées}
Poser et calculer $20-7{,}35$, puis $8{,}4-2{,}76$.
\end{exercice}\end{blocLivre}
\begin{blocLivre}{ex-3-3}\begin{exercice}[1]{Par 10, 100 ou 1 000}
Calculer : $4{,}56\times100$ ; $0{,}37\times10$ ; $2{,}5\times1\,000$ ; $37{,}2:100$ ; $5:1\,000$.
\end{exercice}\end{blocLivre}
\begin{blocLivre}{ex-3-4}\begin{exercice}[1]{Par 0,1 ; 0,01 ; 0,001}
Calculer : $45\times0{,}1$ ; $6{,}2\times0{,}01$ ; $300\times0{,}001$ ; $7\times0{,}01$.
\end{exercice}\end{blocLivre}
\begin{blocLivre}{ex-3-5}\begin{exercice}[1]{Premiers produits}
Calculer $2{,}3\times1{,}4$ ; $0{,}5\times0{,}8$ ; $1{,}2\times3$.
\end{exercice}\end{blocLivre}

\rubrique{Entraînement}
\begin{blocLivre}{ex-3-6}\begin{exercice}[2]{Avec un rectangle}
Le rectangle mesure $3{,}2$ cm sur $2{,}5$ cm. Il est découpé en deux parties. Calculer l’aire de chaque partie, puis en déduire $3{,}2\times2{,}5$.
\label{fig-ex-3-6}\figureGeo{cal-6}
\end{exercice}\end{blocLivre}
\begin{blocLivre}{ex-3-7}\begin{exercice}[2]{Placer la virgule}
On sait que $147\times23=3\,381$. En déduire $14{,}7\times2{,}3$ ; $1{,}47\times23$ ; $0{,}147\times0{,}23$ ; $147\times0{,}023$.
\end{exercice}\end{blocLivre}
\begin{blocLivre}{ex-3-8}\begin{exercice}[2]{Ordre de grandeur}
Sans poser, choisir le bon résultat. $19{,}8\times5{,}1$ : $10{,}098$ ; $100{,}98$ ; $1\,009{,}8$. Puis $0{,}49\times61$ : $2{,}989$ ; $29{,}89$ ; $298{,}9$.
\end{exercice}\end{blocLivre}
\begin{blocLivre}{ex-3-9}\begin{exercice}[2]{Calculer astucieusement}
Calculer de tête en changeant l’ordre des facteurs : $2{,}5\times7{,}3\times4$ ; $0{,}5\times36{,}8\times2$ ; $1{,}25\times3{,}1\times8$.
\end{exercice}\end{blocLivre}
\begin{blocLivre}{ex-3-10}\begin{exercice}[2]{Plus petit ou plus grand ?}
Sans calculer, dire si chaque résultat est plus petit ou plus grand que 24, puis calculer : $24\times0{,}9$ ; $24\times1{,}2$ ; $24\times0{,}05$.
\end{exercice}\end{blocLivre}

\rubrique{Problèmes}
\begin{blocLivre}{ex-3-11}\begin{exercice}[3]{Les courses}
Zoé achète 3 kg de pommes à $2{,}45$ € le kilogramme et $1{,}5$ kg de raisin à $4{,}80$ € le kilogramme. Elle paie avec un billet de 20 €. Combien lui rend-on ?
\end{exercice}\end{blocLivre}
\begin{blocLivre}{ex-3-12}\begin{exercice}[3]{Le terrain}
Un terrain rectangulaire mesure $24{,}5$ m sur $12{,}8$ m. Estimer son aire, puis la calculer exactement.
\end{exercice}\end{blocLivre}
\begin{blocLivre}{ex-3-13}\begin{exercice}[3]{Le plein}
Une voiture consomme $5{,}8$ L de carburant pour 100 km. Quelle quantité consomme-t-elle pour 350 km ? Le carburant coûte $1{,}80$ € le litre : quel est le coût du trajet ?
\end{exercice}\end{blocLivre}
\begin{blocLivre}{ex-3-14}\begin{exercice}[3]{Le ruban}
Un rouleau contient 25 m de ruban. On en coupe 12 morceaux de $1{,}85$ m. Quelle longueur reste-t-il ?
\end{exercice}\end{blocLivre}
\begin{blocLivre}{ex-3-15}\begin{exercice}[3]{L’erreur de virgule}
Nina écrit : $3{,}6\times0{,}5=18$. Expliquer pourquoi c’est impossible, puis corriger.
\end{exercice}\end{blocLivre}

\clearpage
